% Compiled with XeLaTeX (needed for fontspec / real system fonts).
% latexmk -xelatex main.tex

\usepackage{fontspec}
\usepackage[english]{babel}
\usepackage{microtype}

% ---- Type: one warm, well-built serif family carries the whole book.
% Static Lora cuts (OFL, from the Lora project), loaded by explicit path
% so the build is machine-independent and bold is a real bold. ----
\setmainfont{Lora}[
  Path = fonts/,
  UprightFont = Lora-Regular.ttf,
  BoldFont = Lora-Bold.ttf,
  ItalicFont = Lora-Italic.ttf,
  BoldItalicFont = Lora-BoldItalic.ttf,
]
\setmonofont{Latin Modern Mono}[Scale=0.92]
\let\displayfont\normalfont

% ---- Page: a real book trim with a centered text block and equal
% margins. Sidenotes are set as styled notes at the foot of the page
% (see \sidenote in boxes.tex), so no page carries an empty column. ----
\usepackage[paperwidth=7.5in, paperheight=9.5in,
            top=0.9in, bottom=1.0in, hmargin=1.0in,
            headheight=15pt, headsep=18pt]{geometry}

% Kept so full-width elements and older sources still compile; with
% symmetric margins "full width" is simply the text width.
\newlength{\fullw}
\setlength{\fullw}{\textwidth}
\newlength{\fullextra}
\setlength{\fullextra}{0pt}

\usepackage{xcolor}
\usepackage[most]{tcolorbox}
\usepackage{listings}
\usepackage{enumitem}
\usepackage[explicit]{titlesec}
\usepackage{fancyhdr}
\usepackage{graphicx}
\usepackage{amsmath,amssymb}
\usepackage[bottom,hang,flushmargin]{footmisc}
\usepackage{booktabs}
\usepackage{array}
\usepackage{changepage}
\usepackage{needspace}
\usepackage{lettrine}
\usepackage[font=small,labelfont={bf,color=accent},labelsep=colon,
            justification=raggedright,singlelinecheck=false]{caption}
\usepackage{tikz}
\usetikzlibrary{positioning,arrows.meta,calc,fit,backgrounds,
                decorations.pathreplacing,shapes.geometric,shapes.misc}

% ---- Palette. One signature color (accent, a warm terracotta) leads;
% one cool partner (cool, a deep teal) exists so diagrams can show two
% sides of a boundary (user vs. kernel, client vs. server, ours vs.
% theirs) without inventing colors per figure. Everything else is ink
% and paper. ----
\definecolor{ink}{HTML}{2A2825}      % body text, warmer than pure black
\definecolor{accent}{HTML}{A34A1F}   % signature: numbers, labels, rules
\definecolor{cool}{HTML}{25646A}     % diagram partner color, trace boxes
\definecolor{gold}{HTML}{D9A441}     % tiny highlights only (the cursor, a marker)
\definecolor{paper}{HTML}{FBF8F2}    % soft panel fill
\definecolor{codeink}{HTML}{2F3B45}  % terminal text
\definecolor{codebg}{HTML}{F4F1EA}   % terminal background, warm paper tone
\definecolor{warn}{HTML}{8C3B12}     % breakage callouts only

\color{ink}

\usepackage[colorlinks=true, linkcolor=accent!70!black,
            citecolor=accent!70!black, urlcolor=cool,
            bookmarksnumbered=true, linktoc=page]{hyperref}
